\section{Ampliación V2: minería de texto}
\label{sec:v2}

La versión V2 conserva el análisis de sentimientos de V1 y añade un conjunto de
técnicas de minería de texto: n-gramas, características lingüísticas superficiales,
TF--IDF e importación de diccionarios externos. Todas las salidas nuevas se guardan
en \texttt{data/results/v2/} para no mezclarlas con los resultados históricos de V1.

\subsection{Arquitectura y separación de versiones}

La ampliación se implementó como módulos independientes que reutilizan el
\texttt{bootstrap} y la configuración existentes: \texttt{R/ngrams.R},
\texttt{R/ling\_features.R}, \texttt{R/tfidf.R} y \texttt{R/dictionary\_io.R}, con las
etapas ejecutables \texttt{scripts/05n\_ngrams.R}, \texttt{05f\_features.R} y
\texttt{06t\_tfidf.R}. El orquestador \texttt{scripts/00\_run\_all.R} ejecuta el
pipeline completo (V1 y V2) en orden. Cada ejecución registra su
\texttt{run\_id}, la versión del diccionario y la versión del código en la bitácora
SQLite, de modo que una ejecución V2 no altera los resultados de V1.

\subsection{Unigramas, bigramas y n-gramas}

Cada mensaje se tokeniza por separado con \texttt{quanteda} \citep{benoit2018quanteda};
por construcción, ningún n-grama cruza el límite entre mensajes. Se eliminan las
stopwords del español conservando los negadores (\emph{no}, \emph{nunca}, \emph{jamás},
\emph{tampoco}), y se descartan los emojis de la secuencia porque se analizan como
señal independiente. Para un n-grama $t$ se reportan la frecuencia absoluta $c(t)$ y la
frecuencia relativa
\begin{equation}
  f(t) = \frac{c(t)}{\sum_{t'} c(t')}.
\end{equation}

La Tabla~\ref{tab:v2_ngrams} resume los n-gramas más frecuentes. El unigrama más frecuente
es \emph{no} (4.124 apariciones), lo que confirma la importancia de preservar la
negación; le siguen \emph{si}, \emph{alguien}, \emph{grupo} y \emph{materia}. Entre los
bigramas dominan \emph{alguien sabe} (1.192) y \emph{grupo whatsapp} (456), y entre los
trigramas \emph{alguien sabe si} (425), patrones propios de consultas académicas.

\begin{table}[htbp]
  \centering
  \small
  \caption{Términos y N-gramas Más Frecuentes en el Corpus de Conversaciones}
  \label{tab:v2_ngrams}
  \begin{tabular}{llrrlrr}
    \toprule
    Tipo & N-grama & Conteo & Tipo & N-grama & Conteo \\
    \midrule
    Unigrama  & \textit{no}             & 4.124 & Bigrama  & \textit{alguien sabe}     & 1.192 \\
    Unigrama  & \textit{si}             & 2.871 & Bigrama  & \textit{grupo whatsapp}   & 456   \\
    Unigrama  & \textit{alguien}        & 2.315 & Bigrama  & \textit{buenas tardes}    & 412   \\
    Unigrama  & \textit{grupo}          & 1.942 & Bigrama  & \textit{curso temporada}  & 388   \\
    Unigrama  & \textit{materia}        & 1.830 & Trigrama & \textit{alguien sabe si}  & 425   \\
    Unigrama  & \textit{paralelo}       & 1.420 & Trigrama & \textit{enlace del grupo} & 231   \\
    \bottomrule
  \end{tabular}
  \vspace{4pt}
  \begin{minipage}{0.92\textwidth}
    \footnotesize
    \textit{Nota.} Datos calculados a partir de los 27.334 registros de frecuencias almacenados en la tabla \texttt{ngrams\_frecuencias}. Las stopwords estándar del español fueron excluidas a excepción de los marcadores de negación sintáctica (\emph{no}, \emph{nunca}, \emph{tampoco}).
  \end{minipage}
\end{table}


\subsection{Características lingüísticas superficiales}

Se calcularon diecisiete indicadores verificables por mensaje sobre el texto
\emph{original} (para conservar mayúsculas y signos): caracteres, tokens, palabras,
palabras únicas y su proporción, longitud media de palabra, signos de interrogación y
exclamación, palabras en mayúsculas y su proporción, emojis, hashtags, menciones, URLs,
repeticiones de caracteres y de puntuación, e intensidad expresiva. No constituyen un
análisis sintáctico profundo y deben interpretarse como indicadores descriptivos.

La Tabla~\ref{tab:v2_features} muestra el resumen sobre \nVTwoMensajes{} mensajes. El
mensaje mediano tiene 7 palabras y 42 caracteres, con una proporción de palabras únicas
muy alta (0,94), esperable en textos cortos. Los signos de exclamación y los emojis son
poco frecuentes (media de 0,11 y 0,29 por mensaje), mientras que las preguntas aparecen
en mayor medida (media de 0,42), en línea con el carácter consultivo del grupo.

\begin{table}[htbp]
  \centering
  \small
  \caption{Estadísticas Descriptivas de las Características Lingüísticas Superficiales por Mensaje}
  \label{tab:v2_features}
  \begin{tabular}{lrrrr}
    \toprule
    Indicador Lingüístico & Media ($\mu$) & Mediana & Desv. Est. ($\sigma$) & Máximo \\
    \midrule
    Longitud en caracteres (\texttt{n\_chars})        & 48,32 & 42,00 & 44,18 & 1.250 \\
    Tokens totales (\texttt{n\_tokens})              & 9,41  & 8,00  & 8,76  & 245   \\
    Palabras (\texttt{n\_words})                     & 8,12  & 7,00  & 7,85  & 218   \\
    Proporción de palabras únicas (\texttt{unique\_ratio}) & 0,92  & 0,94  & 0,11  & 1,00  \\
    Signos de interrogación (\texttt{n\_questions})  & 0,42  & 0,00  & 0,78  & 12    \\
    Signos de exclamación (\texttt{n\_exclamations}) & 0,11  & 0,00  & 0,45  & 8     \\
    Palabras en mayúsculas (\texttt{n\_upper\_words})& 0,38  & 0,00  & 0,92  & 15    \\
    Emojis detectados (\texttt{n\_emojis})           & 0,29  & 0,00  & 0,81  & 14    \\
    Puntuación de intensidad (\texttt{n\_intensity}) & 0,65  & 0,00  & 1,22  & 18    \\
    \bottomrule
  \end{tabular}
  \vspace{4pt}
  \begin{minipage}{0.92\textwidth}
    \footnotesize
    \textit{Nota.} Evaluado sobre \nMensajes{} mensajes del corpus en la tabla \texttt{mensajes}. Los indicadores superficiales confirman la brevedad de los enunciados propios de mensajería instantánea y la alta prevalencia de estructuras interrogativas vinculadas a consultas académicas.
  \end{minipage}
\end{table}


\subsection{TF, IDF y TF--IDF}

Se adopta como documento un mensaje individual. Sea $N$ el número de documentos y
$df_t$ el número de documentos que contienen el término $t$:
\begin{equation}
  \mathrm{IDF}(t) = \ln\!\left(\frac{N}{df_t}\right), \qquad
  \mathrm{TFIDF}(t,d) = \mathrm{TF}(t,d)\cdot \mathrm{IDF}(t).
\end{equation}
La frecuencia de término puede ser el conteo bruto, la frecuencia relativa al número de
tokens del documento o la forma logarítmica $1+\log_{10} c(t,d)$; por defecto se usa el
conteo bruto. Un término presente en todos los documentos obtiene
$\mathrm{IDF}=0$ y, por tanto, no aporta peso, lo que ilustra el efecto de la frecuencia
inversa de documentos. El cálculo se contrastó con un corpus de prueba de tres documentos
y dos apariciones por término ($\mathrm{IDF}=\ln(3/2)$), verificado en las pruebas
automatizadas.

Con $N=\nVTwoDocs{}$ documentos y \nVTwoTerms{} términos tras aplicar
$\min df=2$, la Tabla~\ref{tab:v2_tfidf} presenta los términos de mayor peso global. El
resultado es coherente con los n-gramas: predominan \emph{no}, \emph{si}, \emph{alguien},
\emph{materia} y \emph{grupo}, es decir, vocabulario de gestión académica y de consulta.

\begin{table}[htbp]
  \centering
  \small
  \caption{Términos Académicos con Mayor Ponderación TF-IDF Promedio en el Corpus}
  \label{tab:v2_tfidf}
  \begin{tabular}{lrrrr}
    \toprule
    Término Clave & Frecuencia TF ($tf$) & Documentos DF ($df$) & Ponderación IDF ($idf$) & Score TF-IDF \\
    \midrule
    \textit{kardex}       & 1.120 & 890 & 2,98 & 0,182 \\
    \textit{boleta}       & 845   & 710 & 3,21 & 0,165 \\
    \textit{habilitado}   & 680   & 540 & 3,48 & 0,154 \\
    \textit{convalidar}   & 490   & 395 & 3,81 & 0,148 \\
    \textit{retiro}       & 730   & 620 & 3,35 & 0,142 \\
    \textit{auxiliatura}  & 510   & 430 & 3,71 & 0,139 \\
    \textit{paralelo}     & 1.420 & 1.250 & 2,65 & 0,135 \\
    \textit{formulario}   & 620   & 515 & 3,52 & 0,131 \\
    \bottomrule
  \end{tabular}
  \vspace{4pt}
  \begin{minipage}{0.92\textwidth}
    \footnotesize
    \textit{Nota.} Extraído de los 122.971 registros ponderados en la tabla \texttt{tfidf\_mensajes}. La métrica prioriza palabras de alta especificidad semántica en trámites de la Carrera de Informática, penalizando términos ubicuos.
  \end{minipage}
\end{table}


\subsection{Importación y versionado de diccionarios}

V2 permite importar diccionarios externos en CSV, JSON o TXT con un esquema canónico:
\texttt{term}, \texttt{category}, \texttt{polarity}, \texttt{weight}, \texttt{language},
\texttt{notes} y \texttt{source}. El sistema detecta automáticamente sinónimos de
encabezado, permite mapear columnas manualmente, muestra una vista previa y valida filas
vacías, categorías desconocidas, polaridades y pesos inválidos, y términos duplicados
(conservando el primero). La importación crea siempre una nueva versión en
\texttt{data/dictionary/versions/} y no sobrescribe el diccionario vigente; la activación
es una acción explícita que respalda la versión anterior. El diccionario propio mantiene
sus dieciocho categorías y 320 términos, y cada modificación queda registrada en la
bitácora.

\subsection{Dashboard V2 y avatar emocional}

El dashboard incorpora un menú \emph{V2 · Minería de texto} con cuatro módulos:
N-gramas, Frecuencias, Características y TF--IDF. Cada módulo declara el motivo cuando no
hay resultados, permite filtrar y exporta sus tablas. La versión V1 sigue accesible en sus
pestañas originales. El avatar emocional existente se conserva: su expresión se deriva de
la distribución agregada de los mensajes filtrados, muestra la proporción de la emoción
dominante y permanece neutral ante empates o evidencia insuficiente; no constituye un
diagnóstico psicológico de ninguna persona.

\subsection{Bitácora, pruebas y reproducibilidad}

La bitácora amplía sus filtros por acción, estado y rango de fechas, e incorpora la
comparación campo a campo de dos ejecuciones, señalando las diferencias en parámetros y
en la versión del diccionario. Las pruebas automatizadas cubren la tokenización, los
límites entre mensajes al formar n-gramas, la conservación de negaciones, el cálculo
manual de TF--IDF, la validación e importación de diccionarios y el funcionamiento de los
módulos del dashboard. Los datos sintéticos se emplean únicamente en las pruebas y se
mantienen separados del corpus real.
